\documentclass[11pt]{article}
\usepackage[T1]{fontenc}
\usepackage{lmodern,microtype}
\usepackage[a4paper,margin=27mm]{geometry}
\usepackage{amsmath,amssymb,amsthm,mathtools,booktabs}
\usepackage{xcolor}
\setlength{\emergencystretch}{2em}
\usepackage[colorlinks=true,linkcolor=blue!50!black,citecolor=blue!50!black,urlcolor=blue!50!black]{hyperref}
\newtheorem{theorem}{Theorem}[section]
\newtheorem{proposition}[theorem]{Proposition}
\newtheorem{lemma}[theorem]{Lemma}
\newtheorem{corollary}[theorem]{Corollary}
\theoremstyle{remark}\newtheorem{remark}[theorem]{Remark}
\newcommand{\E}{\mathbb E}
\newcommand{\Prob}{\mathbb P}
\newcommand{\R}{\mathbb R}
\newcommand{\ii}{\mathrm i}
\newcommand{\dd}{\mathrm d}
\newcommand{\asymerror}{O\!\left(\frac1{u\log u}\right)}
\title{A sharp small-deviation law for random sums\newline built from zeta and Dirichlet zeros}
\author{Research note prepared with ChatGPT\\for Vladimir Reshetnikov}
\date{26 November 2026}
\begin{document}
\maketitle

\begin{abstract}
Assume the Riemann hypothesis and sum independent unit-mean exponential
random variables divided by the squares of the positive zero ordinates.
The resulting random variable has a classical reciprocal-$\xi$ Laplace
transform. We derive an explicit small-deviation asymptotic, including its
constant and first relative correction. The saddle is expressed with the
negative branch of Lambert's $W$ function. A common theorem covers every
positive gamma-convolution shape and, under Dirichlet GRH, every primitive
real nonprincipal Dirichlet character. The conductor, parity, central
vanishing order, and first nonzero central coefficient all occur explicitly.
The proof includes a uniform first-order density approximation under
exponential tilting, sufficient to justify the relative error. The
probabilistic representation and general saddlepoint method are established
ideas; the proposed contribution is this explicit evaluation and extension.
Literature priority for the precise formulas has not been certified.
\end{abstract}

\section{The zeta result in explicit form}

Write
\begin{equation}\label{eq:xi}
 \xi(s)=\frac12s(s-1)\pi^{-s/2}\Gamma(s/2)\zeta(s),
 \qquad c_0=\xi(1/2)>0.
\end{equation}
Under RH, list the positive ordinates of the nontrivial zeros as
$\gamma_1,\gamma_2,\ldots$, repeated according to multiplicity. No simplicity
assumption is made. Let $E_j$ be independent exponential random variables of
mean $1$ and define
\[
 X=\sum_{j\geq1}\frac{E_j}{\gamma_j^2}.
\]
This sum is finite almost surely because its expectation is finite.

\begin{theorem}[Explicit zeta small deviations]\label{thm:zeta}
Assume RH. For sufficiently small $x>0$, let
\begin{equation}\label{eq:zeta-saddle}
 v=-W_{-1}(-8\pi x),\qquad u=\frac{v}{4x}.
\end{equation}
Equivalently, $v>1$ is the solution of $ve^{-v}=8\pi x$ and
$v=\log(u/(2\pi))$. Put
\begin{equation}\label{eq:zeta-A}
 A(x)=\frac{2^{5/4}c_0}{\pi^{3/4}}
       \frac{u^{-9/4}}{\sqrt{v-1}}
       \exp\!\left(-\frac{u(v-2)}4\right).
\end{equation}
Then $X$ has a smooth density $f$ on $(0,\infty)$ and, as $x\downarrow0$,
\begin{align}
 \Prob(X\leq x)&=A(x)
       \left(1+\frac1{48u}+\asymerror\right),\label{eq:zeta-cdf}\\
 f(x)&=u^2A(x)
       \left(1+\frac1{48u}+\asymerror\right).\label{eq:zeta-density}
\end{align}
In particular,
\begin{equation}\label{eq:log-leading}
 \lim_{x\downarrow0}
 \frac{x\log\Prob(X\leq x)}{\log^2(1/x)}=-\frac1{16}.
\end{equation}
\end{theorem}

The explicit expression in \eqref{eq:zeta-A} is a relative asymptotic, not
merely a logarithmic estimate: it fixes the multiplicative constant and the
power of $u$. In terms of $x$ and $v$ alone, its exponent is
$-v(v-2)/(16x)$. Replacing $v$ just by $\log(1/x)$ gives the correct limit
\eqref{eq:log-leading}, but not the relative asymptotic.

The theorem concerns a random variable constructed from all the positive
zero ordinates. It does not assert random behavior of the zeros themselves.
It remains valid with repeated zeros. Its proof below does not use pair
correlation, linear independence of ordinates, or central nonvanishing
results for Dirichlet $L$-functions.

\section{An abstract evaluation theorem}

We give hypotheses that isolate the analytic information used. Let
$\gamma_j>0$ be a multiset with
\begin{equation}\label{eq:count}
 N(T)=\#\{j:\gamma_j\leq T\}
       \sim\frac{T\log T}{2\pi}.
\end{equation}
Suppose $c,D,q>0$, $\beta\in\R$, and
\begin{equation}\label{eq:product}
 H(u)=c\prod_{j\geq1}\left(1+\frac{u^2}{\gamma_j^2}\right).
\end{equation}
The product converges locally uniformly because
$\sum_j\gamma_j^{-2}<\infty$. Assume that, as $u\to+\infty$,
\begin{equation}\label{eq:h-asymptotic}
 \log\frac{H(u)}c
 =\frac u2\left(\log\frac{qu}{2\pi}-1\right)
  +\beta\log u+\log\frac Dc-\frac1{48u}+R(u),
\end{equation}
where $R^{(k)}(u)=O(u^{-2-k})$ for $k=0,1,2$.

For a fixed shape parameter $\theta>0$, take independent gamma random
variables $G_j$ with density
\[
 \frac{y^{\theta-1}e^{-y}}{\Gamma(\theta)}\,\mathbf1_{y>0},
\]
and set
\begin{equation}\label{eq:clock}
 X_\theta=\sum_{j\geq1}\frac{G_j}{\gamma_j^2}.
\end{equation}
Here $\theta=1$ gives the exponential case.

\begin{theorem}[Gamma-sum evaluation]\label{thm:abstract}
Under \eqref{eq:count}--\eqref{eq:h-asymptotic}, let
\begin{equation}\label{eq:general-saddle}
 v=-W_{-1}\!\left(-\frac{8\pi x}{\theta q}\right),
 \qquad u=\frac{\theta v}{4x}.
\end{equation}
Thus $v=\log(qu/(2\pi))$ and $v>1$ for sufficiently small $x>0$.
Define
\begin{equation}\label{eq:general-A}
 A_\theta(x)=
 \frac{2(c/D)^\theta}{\sqrt{\pi\theta}}
 \frac{u^{-\theta\beta-1/2}}{\sqrt{v-1}}
 \exp\!\left(-\frac{\theta u(v-2)}4\right).
\end{equation}
The density $f_\theta$ of $X_\theta$ satisfies
\begin{align}
 \Prob(X_\theta\leq x)
 &=A_\theta(x)\left(1+\frac{\theta}{48u}
                         +\asymerror\right),\label{eq:general-cdf}\\
 f_\theta(x)
 &=u^2A_\theta(x)\left(1+\frac{\theta}{48u}
                         +\asymerror\right).\label{eq:general-density}
\end{align}
The constants implicit in $O(\cdot)$ may depend on the fixed data and
$\theta$. No uniformity over growing conductors or varying $\theta$ is
asserted. In particular,
\[
 \lim_{x\downarrow0}
 \frac{x\log\Prob(X_\theta\leq x)}{\log^2(1/x)}
 =-\frac{\theta^2}{16}.
\]
\end{theorem}

\section{Exponential tilting and a quantitative local limit}

Put $h(u)=\log(H(u)/c)$ and
\begin{equation}\label{eq:K}
 K(t)=\theta h(\sqrt t)
     =\theta\sum_j\log(1+t/\gamma_j^2),\qquad t\geq0.
\end{equation}
Independence and bounded convergence give the exact transform
\begin{equation}\label{eq:laplace}
 \E e^{-tX_\theta}=e^{-K(t)}
     =\left(\frac c{H(\sqrt t)}\right)^\theta.
\end{equation}
For $t>0$, define the tilted probability measure by
\[
 \frac{\dd\Prob_t}{\dd\Prob}=e^{-tX_\theta+K(t)}.
\]
Under $\Prob_t$, the summands in \eqref{eq:clock} remain independent gamma
variables, with shape $\theta$ and rates $\gamma_j^2+t$. Their total mean
and variance are
\begin{equation}\label{eq:mean-var}
 m(t)=K'(t)=\theta\sum_j(\gamma_j^2+t)^{-1},\qquad
 \sigma^2(t)=-K''(t)=\theta\sum_j(\gamma_j^2+t)^{-2}.
\end{equation}
These identities follow first for finitely many summands and then by
dominated convergence. The tilted joint transforms of every finite set of
summands factor; hence the stated independence also holds for the infinite
sum.

Differentiating \eqref{eq:h-asymptotic}, with $t=u^2$ and
$\ell=\log(qu/(2\pi))$, gives
\begin{align}
 m(u^2)&=\theta\left(\frac{\ell}{4u}
                  +\frac{\beta}{2u^2}
                  +\frac1{96u^3}+O(u^{-4})\right),\label{eq:mexp}\\
 \sigma^2(u^2)&=\theta\left(\frac{\ell-1}{8u^3}
                  +\frac{\beta}{2u^4}+O(u^{-5})\right).
                  \label{eq:sexp}
\end{align}
In particular $\sigma^2(u^2)\asymp(\log u)/u^3$.

\begin{lemma}[Uniform density approximation]\label{lem:edgeworth}
Let $p_t$ be the density under $\Prob_t$ of
\[
 Y_t=\frac{X_\theta-m(t)}{\sigma(t)}.
\]
Write $\varphi(y)=(2\pi)^{-1/2}e^{-y^2/2}$,
$H_3(y)=y^3-3y$, and
\[
 \kappa_3(t)=
  2\theta\sum_j\bigl[\sigma(t)(\gamma_j^2+t)\bigr]^{-3}.
\]
As $t=u^2\to\infty$, with $n=u\log u$,
\begin{equation}\label{eq:edgeworth}
 \sup_{y\in\R}\left|
  p_t(y)-\varphi(y)-\frac{\kappa_3(t)}6H_3(y)\varphi(y)
 \right|=O(n^{-1}),\qquad \kappa_3(t)=O(n^{-1/2}).
\end{equation}
\end{lemma}

\begin{proof}
Let $b_j=[\sigma(t)(\gamma_j^2+t)]^{-1}$. Then
$\theta\sum_jb_j^2=1$ and
\[
 \max_jb_j\leq\frac1{t\sigma(t)}=O(n^{-1/2}).
\]
The centered characteristic function is
\[
 \phi_t(w)=\exp\!\left(-\frac{\ii wm(t)}{\sigma(t)}\right)
            \prod_j(1-\ii wb_j)^{-\theta}.
\]
The count \eqref{eq:count} shows that there are at least $a_1 n$ ordinates
in $[u,2u]$, for a constant $a_1>0$ and all sufficiently large $u$.
For these ordinates $b_j^2\geq a_2/n$ by \eqref{eq:sexp}. Consequently
there are fixed constants $a_3,a_4>0$ such that
\begin{equation}\label{eq:fourier-tail}
 |\phi_t(w)|\leq(1+a_3w^2/n)^{-a_4 n}
 \quad(w\in\R).
\end{equation}

For $|w|\leq n^{1/12}$, expand the logarithm of the characteristic
function. The linear term cancels the centering term, and
\[
 \log\phi_t(w)=-\frac{w^2}{2}
                -\frac{\ii\kappa_3(t)w^3}{6}
                +O(w^4/n),
\]
because
$\theta\sum_jb_j^4\leq(\max_jb_j)^2=O(n^{-1})$.
The constants in this estimate are uniform on that interval. Expanding
the exponential one more time therefore yields
\begin{equation}\label{eq:fourier-local}
 \phi_t(w)=e^{-w^2/2}
          \left(1-\frac{\ii\kappa_3(t)w^3}{6}\right)
          +e^{-w^2/2}O\!\left(\frac{w^4+w^6}{n}\right).
\end{equation}
For sufficiently large $n$, the small exponential of the remainder can
be absorbed into a fixed constant on this interval.

For completeness, \eqref{eq:fourier-tail} controls the complementary
integral, not just fixed frequencies. If
$n^{1/12}\leq |w|\leq\sqrt n$, the bound is at most $e^{-a_5w^2}$.
If $|w|=\sqrt n z\geq\sqrt n$, then
\[
 1+a_3z^2\geq(1+a_3)z^{2a_3/(1+a_3)}\quad(z\geq1),
\]
as follows by differentiating the logarithms. Thus the integral of the
bound on $|w|\geq\sqrt n$ is $O(e^{-a_6 n})$. The integral on the
intermediate range is $O(e^{-a_7n^{1/6}})$. The Gaussian comparison has
the same negligible tails. Integrating \eqref{eq:fourier-local} proves
\[
 \int_\R\left|
  \phi_t(w)-e^{-w^2/2}(1-\ii\kappa_3(t)w^3/6)
 \right|\dd w=O(n^{-1}).
\]
Fourier inversion, and
$-\varphi'''(y)=H_3(y)\varphi(y)$, prove the uniform density estimate.
The bound on $\kappa_3$ follows from
$\theta\sum_jb_j^3\leq\max_jb_j$.
\end{proof}

\begin{lemma}[Relative saddlepoint formulas]\label{lem:saddle}
For $0<x<\E X_\theta$, let $t_*=t_*(x)$ be the unique solution of
$m(t_*)=x$, and set $u_*=\sqrt{t_*}$. Then, as $x\downarrow0$,
\begin{align}
 \Prob(X_\theta\leq x)
 &=\frac{e^{t_*x-K(t_*)}}{
          t_*\sqrt{2\pi\sigma^2(t_*)}}
     \left(1+O\!\left(\frac1{u_*\log u_*}\right)\right),
     \label{eq:implicit-cdf}\\
 f_\theta(x)
 &=\frac{e^{t_*x-K(t_*)}}{
          \sqrt{2\pi\sigma^2(t_*)}}
     \left(1+O\!\left(\frac1{u_*\log u_*}\right)\right).
     \label{eq:implicit-density}
\end{align}
\end{lemma}

\begin{proof}
The mean $m(t)$ is continuous and strictly decreasing from
$\E X_\theta$ to $0$, by \eqref{eq:mean-var}. This proves existence,
uniqueness, and $t_*\to\infty$.
For $t=t_*$ put $a=t\sigma(t)$. The exact change of probability measure is
\begin{align*}
 \Prob(X_\theta\leq x)
 &=e^{tx-K(t)}\E_t\left[e^{t(X_\theta-x)}
                              \mathbf1_{X_\theta\leq x}\right]\\
 &=\frac{e^{tx-K(t)}}a
    \int_0^\infty e^{-z}p_t(-z/a)\dd z.
\end{align*}
By \eqref{eq:sexp}, $a\asymp\sqrt n$. The elementary bounds
\[
 |\varphi(y)-\varphi(0)|\leq C y^2,
 \qquad |H_3(y)\varphi(y)|\leq C|y|
\]
hold for all real $y$. Lemma~\ref{lem:edgeworth} therefore gives
\[
 \int_0^\infty e^{-z}p_t(-z/a)\dd z
 =\frac1{\sqrt{2\pi}}
   +O(a^{-2}+|\kappa_3(t)|a^{-1}+n^{-1})
 =\frac1{\sqrt{2\pi}}+O(n^{-1}).
\]
This proves \eqref{eq:implicit-cdf}. The exact density identity is
$f_\theta(x)=e^{tx-K(t)}p_t(0)/\sigma(t)$; since $H_3(0)=0$, the
same lemma proves \eqref{eq:implicit-density}.

The required densities exist and are smooth. Indeed, the characteristic
function of the original sum, or of any tilt, has arbitrarily high
polynomial decay: retain as many factors in its product as needed, and
bound the remaining factors by $1$. Fourier inversion then supplies all
fixed-order density derivatives.
\end{proof}

\section{Evaluating the saddle and the first correction}

\begin{proof}[Proof of Theorem~\ref{thm:abstract}]
Let $u,v$ be the explicit values in \eqref{eq:general-saddle}. They solve
\begin{equation}\label{eq:leading-saddle}
 x=\frac{\theta\log(qu/(2\pi))}{4u}=\frac{\theta v}{4u}.
\end{equation}
Let $u_*$ denote the true saddle from Lemma~\ref{lem:saddle}. Formula
\eqref{eq:mexp} and the mean value theorem give
\begin{equation}\label{eq:saddle-shift}
 u_*-u=O((\log u)^{-1}).
\end{equation}
To see the scale explicitly, the mean error at $u$ is
$\theta\beta/(2u^2)+O(u^{-3})$, whereas the derivative of the leading
mean with respect to $u$ is
$-\theta(v-1)/(4u^2)$. These estimates also bracket the true saddle in
an interval of the claimed radius, without assuming its location in
advance.

Define $\Phi(w)=xw^2-\theta h(w)$. At $w=u$,
\[
 \Phi'(u)=-\frac{\theta\beta}{u}+O(u^{-2}),
 \qquad
 \Phi''(w)=O((\log u)/u)
 \quad\text{if }|w-u|=O((\log u)^{-1}).
\]
Taylor's theorem and \eqref{eq:saddle-shift} imply
\begin{equation}\label{eq:action-shift}
 \Phi(u_*)-\Phi(u)=O((u\log u)^{-1}).
\end{equation}
The same shift changes either saddlepoint denominator by a relative
$O((u\log u)^{-1})$. Indeed, \eqref{eq:sexp} gives
\begin{equation}\label{eq:denominator}
 u^2\sqrt{2\pi\sigma^2(u^2)}
 =\frac{\sqrt{\pi\theta}}2\,
   u^{1/2}\sqrt{v-1}
   \left(1+O((u\log u)^{-1})\right),
\end{equation}
and the logarithmic derivatives of the explicit power and square-root
factors are $O(1/u)$.

Finally, \eqref{eq:h-asymptotic} and \eqref{eq:leading-saddle} give
\begin{align*}
 e^{\Phi(u)}
 &=\left(\frac cD\right)^\theta u^{-\theta\beta}
   \exp\!\left(-\frac{\theta u(v-2)}4\right)
   \exp\!\left(\frac{\theta}{48u}+O(u^{-2})\right)\\
 &=\left(\frac cD\right)^\theta u^{-\theta\beta}
   \exp\!\left(-\frac{\theta u(v-2)}4\right)
   \left(1+\frac{\theta}{48u}+O(u^{-2})\right).
\end{align*}
Insert this, \eqref{eq:action-shift}, and \eqref{eq:denominator} into
Lemma~\ref{lem:saddle}. This proves both relative formulas. Since
$v\sim\log(1/x)$ and $u=\theta v/(4x)$, their leading exponent is
$-\theta^2\log^2(1/x)/(16x)$, proving the final limit.
\end{proof}

\section{Verification for zeta and real Dirichlet characters}

\subsection{Zeta}

The paired Hadamard product, RH, and the functional equation give
\[
 \xi(1/2+u)=c_0\prod_{j\geq1}(1+u^2/\gamma_j^2).
\]
The Riemann--von Mangoldt count \cite{titchmarsh} supplies \eqref{eq:count}.
Stirling's expansion \cite{dlmf-gamma}, together with
$\log\zeta(u+1/2)=O(2^{-u})$, gives
\begin{equation}\label{eq:xi-explicit}
 \log\xi(1/2+u)
 =\frac u2\left(\log\frac{u}{2\pi}-1\right)
  +\frac74\log u+\frac14\log\frac\pi2
  -\frac1{48u}+O(u^{-2}).
\end{equation}
The expansion can be differentiated twice with the stated remainder
bounds. One can use the differentiated Stirling expansion and the
absolutely convergent Dirichlet series, or Cauchy estimates in a fixed
sector around the positive axis.

Thus Theorem~\ref{thm:abstract} applies with
\[
 q=1,\qquad \beta=\frac74,\qquad
 D=(\pi/2)^{1/4},\qquad c=c_0.
\]
Taking $\theta=1$ proves Theorem~\ref{thm:zeta}; arbitrary $\theta>0$
also gives the gamma-convolution extension for zeta.

\subsection{Real primitive Dirichlet characters}

Let $\chi$ be a primitive real nonprincipal Dirichlet character of conductor
$q$. Define $a\in\{0,1\}$ by $\chi(-1)=(-1)^a$, and use the completion
\begin{equation}\label{eq:completion}
 \Lambda_\chi(s)=
 \left(\frac q\pi\right)^{(s+a)/2}
 \Gamma\!\left(\frac{s+a}{2}\right)L(s,\chi).
\end{equation}
Let $m$ be its order of vanishing at $s=1/2$ and put
\begin{equation}\label{eq:central}
 c_\chi=\frac{\Lambda_\chi^{(m)}(1/2)}{m!},\quad
 \beta_\chi=\frac{2a-1}{4}-m,\quad
 D_\chi=2\sqrt\pi
          \left(\frac q{2\pi}\right)^{(2a+1)/4}.
\end{equation}
When $m=0$, $c_\chi=\Lambda_\chi(1/2)$. In general these are definitions
using the first nonzero central derivative, so no central nonvanishing
assumption is needed.

\begin{corollary}[Explicit real-Dirichlet law]\label{cor:Dirichlet}
Assume GRH for $L(s,\chi)$. List its positive zero ordinates $\gamma_j$
with multiplicity, excluding central zeros, and define $X_{\chi,\theta}$
as in \eqref{eq:clock}. Then $c_\chi>0$. For
\[
 v=-W_{-1}\!\left(-\frac{8\pi x}{\theta q}\right),\qquad
 u=\frac{\theta v}{4x},
\]
set
\begin{equation}\label{eq:Dirichlet-A}
 A_{\chi,\theta}(x)=
 \frac{2(c_\chi/D_\chi)^\theta}{\sqrt{\pi\theta}}
 \frac{u^{\theta m-\theta(2a-1)/4-1/2}}{\sqrt{v-1}}
 \exp\!\left(-\frac{\theta u(v-2)}4\right).
\end{equation}
The CDF and density satisfy, respectively,
\[
 \Prob(X_{\chi,\theta}\leq x)
 =A_{\chi,\theta}(x)
     \left(1+\frac{\theta}{48u}+\asymerror\right),
\]
\[
 f_{\chi,\theta}(x)
 =u^2A_{\chi,\theta}(x)
     \left(1+\frac{\theta}{48u}+\asymerror\right).
\]
\end{corollary}

\begin{proof}
The completion \eqref{eq:completion} is entire of order one. For real
$\chi$, its functional equation has a real unit root number, so dividing
$\Lambda_\chi(1/2+u)$ by its central power $u^m$ produces an even entire
function
\[
 H_\chi(u)=\Lambda_\chi(1/2+u)/u^m,
 \qquad H_\chi(0)=c_\chi.
\]
The removable value at $0$ is understood. To verify evenness directly,
write the functional equation as
$\Lambda_\chi(1/2+u)=\varepsilon\Lambda_\chi(1/2-u)$.
Its first nonzero Taylor coefficient gives $\varepsilon(-1)^m=1$.
Consequently $H_\chi(u)=H_\chi(-u)$.

Under GRH all its zeros are $\pm\ii\gamma_j$. The paired order-one
Hadamard product gives
\[
 H_\chi(u)=c_\chi\prod_j(1+u^2/\gamma_j^2).
\]
There is no additional linear exponential factor because the function is
even, and no quadratic exponential factor because it has order one.
For real $s>1$, the Euler product makes $L(s,\chi)>0$; GRH implies no
zeros on the real interval $s>1/2$. Continuity therefore gives
$\Lambda_\chi(1/2+u)>0$ for $u>0$, and its leading central coefficient
$c_\chi$ is positive. The usual zero count for a fixed Dirichlet character
is \eqref{eq:count}; removing finitely many central zeros does not affect
it. The standard completion, functional equation, and zero structure are
reviewed in \cite{dlmf-dirichlet,davenport}.

To check the constants, put $b=(2a+1)/4$, so that the gamma argument in
\eqref{eq:completion} is $u/2+b$. Stirling expansion gives
\begin{align*}
 \log H_\chi(u)
 &=\frac u2\left(\log\frac{qu}{2\pi}-1\right)
   +(b-1/2-m)\log u\\
 &\quad+\log(2\sqrt\pi)+b\log\frac q{2\pi}
   +\frac{b^2-b+1/6}{u}+O(u^{-2}).
\end{align*}
For both $b=1/4$ and $b=3/4$,
$b^2-b+1/6=-1/48$. Also $L(u+1/2,\chi)=1+O(2^{-u})$,
with the corresponding derivative estimates. This verifies
\eqref{eq:h-asymptotic} and its differentiated remainders with the data
\eqref{eq:central}. Theorem~\ref{thm:abstract} proves the result.
\end{proof}

\section{What is established and what is proposed as new}

The reciprocal-$\xi$ probability law is not new. Wald couples, reciprocal
transforms, and their connections with RH belong to established work,
including Roynette--Yor \cite{roynette-yor}. General small-deviation
saddlepoint principles also predate this note; Hadamard-factorization
methods for evaluating such laws are developed in
Gao--Hannig--Lee--Torcaso \cite{gao}.

A particularly close predecessor is Sra's 2 November 2026 preprint
\cite{sra}, Appendix C.17, Proposition C.25. It studies the same
reciprocal-$\xi$ density and proves convergence of tilted density
derivatives to Gaussian density derivatives, using Stirling asymptotics,
in order to establish local total positivity. Thus neither a tilted local
Gaussian limit nor the use of Stirling's formula for this law is claimed
as a new idea here.

The candidate contribution of this note is the explicit Lambert-$W$
evaluation of the CDF and density, their complete multiplicative
constants, the isolated relative correction $\theta/(48u)$ with a
smaller error, and the single formula retaining conductor, parity, and
possible central vanishing for every primitive real Dirichlet character.
Targeted searches for the corresponding small-deviation, lower-tail,
Lambert-$W$, and Dirichlet formulas did not locate this precise statement.
That search is not a proof of historical priority. The result should be
described as a proved formula whose novelty requires further specialist
review, rather than as a certified first publication.

The theorem uses RH/Dirichlet GRH only to identify the completed functions'
zeros with positive spectral rates. The cited Lean repository
\cite{grhlean} supplies exactly these zero-location inputs at the scope
stated by that project; its automorphic argument is not reproduced here.
All probability and asymptotic arguments in this note are human-readable
proofs, not Lean formalizations. No result about general automorphic or
Artin $L$-functions is asserted.

\section{Reproducible numerical check}

The companion script \texttt{check\_xi\_small\_deviations.py} evaluates the
zeta transform under the true exponential tilt and then integrates its
standardized Fourier transform numerically. It uses 60 decimal digits and
truncates at standardized frequency 20. The values below compare those
numerical integrals with the leading explicit formulas, using
$x=m(u_*^2)$ for the displayed true saddle $u_*$.

\begin{center}
\begin{tabular}{@{}rrr@{}}
\toprule
$u_*$ & Numerical $\Prob(X\leq x)/A(x)$ & Numerical $f(x)/(u^2A(x))$\\
\midrule
$100$   & $0.9351560720$ & $0.9882302989$\\
$1000$  & $0.9972922246$ & $0.9996171792$\\
$10000$ & $0.9998321209$ & $0.9999782512$\\
\bottomrule
\end{tabular}
\end{center}

These are exploratory floating-point checks, not interval certificates;
the script does not provide a rigorous truncation error. They are not used
in the proof. The first isolated relative term $1/(48u)$ is asymptotically
larger than the remainder, but its small coefficient means that it need
not improve the leading approximation at the modest saddles in this
table. The accompanying JSON preserves additional digits, the input
values, and the precision settings.

\begin{thebibliography}{9}

\bibitem{roynette-yor}
B.~Roynette and M.~Yor,
\emph{Couples de Wald ind\'efiniment divisibles. Exemples li\'es \`a la
fonction gamma d'Euler et \`a la fonction zeta de Riemann},
Annales de l'Institut Fourier \textbf{55} (2005), 1219--1283.
\href{https://doi.org/10.5802/aif.2125}{doi:10.5802/aif.2125}.

\bibitem{gao}
F.~Gao, J.~Hannig, T.-Y.~Lee, and F.~Torcaso,
\emph{Laplace Transforms via Hadamard Factorization},
Electronic Journal of Probability \textbf{8} (2003), paper 13, 1--20.
\url{https://www.maths.tcd.ie/EMIS/journals/EJP-ECP/article/view/151.html}.

\bibitem{sra}
S.~Sra,
\emph{On some positivity properties of the Riemann $\Xi$ function},
preprint, version 3, 2 November 2026; especially Appendix C.17,
Proposition C.25.
\url{https://suvrit.github.io/papers/riem-xi/riem-xi-v3.pdf}.

\bibitem{dlmf-gamma}
NIST Digital Library of Mathematical Functions,
\S5.11, \emph{Asymptotic Expansions}, equations 5.11.1--5.11.2.
\url{https://dlmf.nist.gov/5.11}. Accessed 26 November 2026.

\bibitem{dlmf-dirichlet}
NIST Digital Library of Mathematical Functions,
\S25.15, \emph{Dirichlet $L$-functions}.
\url{https://dlmf.nist.gov/25.15}. Accessed 26 November 2026.

\bibitem{davenport}
H.~Davenport,
\emph{Multiplicative Number Theory}, third edition, revised by H.~L.
Montgomery, Graduate Texts in Mathematics 74, Springer, 2000.

\bibitem{titchmarsh}
E.~C.~Titchmarsh,
\emph{The Theory of the Riemann Zeta-Function}, second edition, revised by
D.~R.~Heath-Brown, Clarendon Press, 1986.

\bibitem{grhlean}
L.~Sernau and contributors,
\emph{grh-lean: GRH for Dirichlet $L$-functions, in Lean 4},
repository and reading guide, commit\
\texttt{72fead102f663d5c7a407964a13dc0cf9d977207},
consulted 26 November 2026.
\url{https://github.com/lsernau/grh-lean}.

\end{thebibliography}
\end{document}
